\begin{proof}[Proof of \cref{obj:lem:synthetic-randomization-kernel}]
\begin{enumerate}
\item Construct the full-data completion. For observational categories \(x\in\{1,\ldots,d\}\) and \(b\in\{0,1\}\), use the following total versions of the quantities in \cref{obj:def:zeng-discrete-ate-class}:
\[
p_x:=P_Z(X=x),\qquad
\pi_x:=
\begin{cases}
P_Z(X=x,B^{\mathrm{obs}}=1)/p_x,&p_x>0,\\
0,&p_x=0,
\end{cases}
\]
and
\[
\mu_{bx}:=
\begin{cases}
\dfrac{P_Z(X=x,B^{\mathrm{obs}}=b,Y=1)}
      {P_Z(X=x,B^{\mathrm{obs}}=b)},
&P_Z(X=x,B^{\mathrm{obs}}=b)>0,\\[6pt]
0,&P_Z(X=x,B^{\mathrm{obs}}=b)=0.
\end{cases}
\]
The two observational arm probabilities sum to \(p_x\). Positivity therefore gives, when \(p_x>0\),
\[
P_Z(X=x,B^{\mathrm{obs}}=1)\ge qp_x>0,\qquad
P_Z(X=x,B^{\mathrm{obs}}=0)\ge qp_x>0.
\]
Thus these means agree with those in \cref{obj:def:zeng-discrete-ate-class} on occupied categories. Each ratio lies in \([0,1]\), since its numerator is the probability of a subset of the event in its denominator.

Define the binary mass functions
\[
\beta^{\mathrm{mark}}_{bx}:=
\begin{cases}
1-\pi_x,&b=0,\\
\pi_x,&b=1,
\end{cases}
\qquad
\beta^{\mathrm{out}}_{bxy}:=
\begin{cases}
1-\mu_{bx},&y=0,\\
\mu_{bx},&y=1,
\end{cases}
\quad
(x\in\{1,\ldots,d\},\ b,y\in\{0,1\}).
\]
Their masses are nonnegative and satisfy
\[
\sum_{b=0}^1\beta^{\mathrm{mark}}_{bx}=1,
\qquad
\sum_{y=0}^1\beta^{\mathrm{out}}_{bxy}=1.
\]
Let \(\mathbb P_{\mathrm{lat}}\) be the probability law on
\(\{0,\ldots,d-1\}\times\{0,1\}^4\), with coordinates
\((X,B^{\mathrm{lat}},Y(0),Y(1),A)\), specified by
\[
\mathbb P_{\mathrm{lat}}
\bigl(X=x-1,B^{\mathrm{lat}}=b,Y(0)=y_0,Y(1)=y_1,A=a\bigr)
=
\frac{p_x}{2}\,
\beta^{\mathrm{mark}}_{bx}
\beta^{\mathrm{out}}_{0xy_0}
\beta^{\mathrm{out}}_{1xy_1},
\]
for \(x\in\{1,\ldots,d\}\) and \(a,b,y_0,y_1\in\{0,1\}\).
The preceding normalization identities and \(\sum_xp_x=1\) show that this is a probability law.

On this probability space set
\[
S(0):=S(1):=0,\qquad
S:=S(A),\qquad
Y:=Y(A),\qquad
R:=\mathbf 1\{B^{\mathrm{lat}}=A\},
\]
and define the full-data law
\[
P_Z^*
:=
\mathcal L_{\mathbb P_{\mathrm{lat}}}
\bigl(X,A,S(0),S(1),Y(0),Y(1),S,Y,R\bigr).
\]
Here \(\mathcal L\) denotes the distribution of the displayed coordinates.
% lean: CausalSmith.Stat.MarRareqLogfrontier.SyntheticCompletion.fullLaw

\item Verify membership in the arrival-floor model. The product factor \(1/2\) in the construction gives
\[
A\perp\!\!\!\perp
\bigl(X,S(0),S(1),Y(0),Y(1)\bigr),
\qquad
P_Z^*(A=1)=\frac12.
\]
Together with \(S=S(A)\) and \(Y=Y(A)\), these establish
\cref{obj:ass:randomized-independence,obj:ass:balanced-randomization,obj:ass:surrogate-consistency,obj:ass:outcome-consistency}.

For \(a,r_{\mathrm{arr}}\in\{0,1\}\) and \(x\in\{1,\ldots,d\}\), define the arrival mass
\[
\beta^{\mathrm{arr}}_{axr_{\mathrm{arr}}}
:=
\begin{cases}
\beta^{\mathrm{mark}}_{ax},&r_{\mathrm{arr}}=1,\\
\beta^{\mathrm{mark}}_{1-a,x},&r_{\mathrm{arr}}=0.
\end{cases}
\]
Fix also \(y\in\{0,1\}\). The event \(R=r_{\mathrm{arr}}\), \(A=a\) fixes the latent mark to \(a\) when \(r_{\mathrm{arr}}=1\), and to \(1-a\) when \(r_{\mathrm{arr}}=0\). Summing over the unselected potential outcome gives
\begin{equation}\label{eq:synthetic-randomization-joint-mass}
P_Z^*(A=a,X=x-1,S=0,R=r_{\mathrm{arr}},Y=y)
=
\frac{p_x}{2}\,
\beta^{\mathrm{arr}}_{axr_{\mathrm{arr}}}
\beta^{\mathrm{out}}_{axy}.
\end{equation}
Indeed, the mass of the unselected potential outcome sums to one. Summing this identity over the binary outcome and arrival coordinates yields
\[
\begin{aligned}
p_{a,x-1,0}(P_Z^*)&=\frac{p_x}{2},\\
P_Z^*(A=a,X=x-1,S=0,R=r_{\mathrm{arr}})
&=\frac{p_x}{2}\beta^{\mathrm{arr}}_{axr_{\mathrm{arr}}},\\
P_Z^*(A=a,X=x-1,S=0,Y=y)
&=\frac{p_x}{2}\beta^{\mathrm{out}}_{axy}.
\end{aligned}
\]
Consequently,
\[
\begin{aligned}
&P_Z^*(A=a,X=x-1,S=0,R=r_{\mathrm{arr}},Y=y)\,
p_{a,x-1,0}(P_Z^*)\\
&\quad=
P_Z^*(A=a,X=x-1,S=0,R=r_{\mathrm{arr}})
P_Z^*(A=a,X=x-1,S=0,Y=y).
\end{aligned}
\]
For \(S=1\), every probability in this identity is zero, since \(S=0\) throughout the construction. This proves \cref{obj:ass:arrival-mar}, including cells of zero mass.

The same calculation gives
\[
P_Z^*(A=a,X=x-1,S=0,R=1)
=\frac{p_x}{2}\beta^{\mathrm{mark}}_{ax},
\qquad
p_{a,x-1,1}(P_Z^*)=0.
\]
Thus an occupied cell has \(S=0\) and \(p_x>0\). By the observational positivity bounds,
\[
\begin{aligned}
\rho_{0,x-1,0}(P_Z^*)&=1-\pi_x\ge q,\\
\rho_{1,x-1,0}(P_Z^*)&=\pi_x\ge q.
\end{aligned}
\]
This establishes \cref{obj:ass:occupied-cell-arrival}. Finally, \(n\ge1\), membership of \(P_Z\) in the observational class supplies \(d\ge1\), and \(0<q<1/2\) implies \(q\le1\). Hence
\[
P_Z^*\in\mathcal M^+_{n,d,q}
\]
by \cref{obj:def:unrestricted-arrival-model-class}.
% lean: CausalSmith.Stat.MarRareqLogfrontier.SyntheticCompletion.fullLaw_randomized
% lean: CausalSmith.Stat.MarRareqLogfrontier.SyntheticCompletion.fullLaw_balanced
% lean: CausalSmith.Stat.MarRareqLogfrontier.SyntheticCompletion.fullLaw_surrogate
% lean: CausalSmith.Stat.MarRareqLogfrontier.SyntheticCompletion.fullLaw_outcome
% lean: CausalSmith.Stat.MarRareqLogfrontier.SyntheticCompletion.fullLaw_mar
% lean: CausalSmith.Stat.MarRareqLogfrontier.SyntheticCompletion.fullLaw_arrival

\item Record the sampling property of the canonical observed product law. Define the observed-record alphabet and its one-record law by
\[
\mathcal O_d:=\{0,\ldots,d-1\}\times\{0,1\}^4,
\qquad
\nu_{\mathrm{obs}}
:=\mathcal L_{P_Z^*}(X,A,S,R,RY).
\]
For \(i\in\{1,\ldots,n\}\), let
\[
\operatorname{pr}_i:\mathcal O_d^n\longrightarrow\mathcal O_d,
\qquad
\operatorname{pr}_i(\mathbf o):=\mathbf o_i.
\]
For arbitrary subsets \(E_i\subseteq\mathcal O_d\), the defining property of the product measure gives
\[
\nu_{\mathrm{obs}}^{\otimes n}
\left(\bigcap_{i=1}^n\operatorname{pr}_i^{-1}(E_i)\right)
=
\prod_{i=1}^n\nu_{\mathrm{obs}}(E_i).
\]
The coordinate projections are measurable on these finite spaces. Taking all but one \(E_i\) to be the entire alphabet also gives the common marginal \(\nu_{\mathrm{obs}}\). Thus this product experiment satisfies \cref{obj:ass:iid-sampling}.
% lean: CausalSmith.Stat.MarRareqLogfrontier.sampleLaw_iid

\item Identify the kernel output law. First, the selected latent record has exactly the observational law:
\[
\mathcal L_{\mathbb P_{\mathrm{lat}}}
\bigl(X+1,B^{\mathrm{lat}},Y(B^{\mathrm{lat}})\bigr)=P_Z.
\]
To verify this identity, sum over \(A\) and the unselected potential outcome. For each \(x\in\{1,\ldots,d\}\) and \(b,y\in\{0,1\}\), the resulting singleton mass is
\[
\mathbb P_{\mathrm{lat}}
\bigl(X=x-1,B^{\mathrm{lat}}=b,Y(b)=y\bigr)
=
p_x\beta^{\mathrm{mark}}_{bx}\beta^{\mathrm{out}}_{bxy}.
\]
When \(p_x>0\),
\[
p_x\beta^{\mathrm{mark}}_{bx}
=P_Z(X=x,B^{\mathrm{obs}}=b).
\]
For \(y=1\), multiplying by \(\mu_{bx}\) gives the corresponding observational singleton probability by its defining ratio. For \(y=0\), subtracting that singleton probability from the arm probability gives the same conclusion:
\[
p_x\beta^{\mathrm{mark}}_{bx}\beta^{\mathrm{out}}_{bxy}
=P_Z(X=x,B^{\mathrm{obs}}=b,Y=y).
\]
When \(p_x=0\), both sides vanish because the singleton event is contained in the null baseline category. Equality of all singleton masses proves the selected-record identity.

Define the fair binary law and one-record kernel by
\[
\nu_{\mathrm{coin}}:=\frac12\delta_0+\frac12\delta_1,
\qquad
\kappa_{\mathrm{syn}}\bigl((x,b,y),\cdot\bigr)
:=
\frac12\delta_{o(x,b,y;0)}
+\frac12\delta_{o(x,b,y;1)}.
\]
This kernel is a measurable probability kernel on finite alphabets and depends solely on the input record. Under the completion, the observed tuple satisfies the pointwise identity
\[
(X,A,S,R,RY)
=
o\bigl(X+1,B^{\mathrm{lat}},Y(B^{\mathrm{lat}});A\bigr).
\]
For a matching assignment, the selected outcome equals \(Y(A)\); for a mismatching assignment, both arrived-outcome coordinates are zero. Moreover, \(A\) is independent of the selected latent record. The preceding distributional identity therefore gives
\[
\nu_{\mathrm{obs}}
=
P_Z\kappa_{\mathrm{syn}}
:=
\int\kappa_{\mathrm{syn}}\bigl((x,b,y),\cdot\bigr)\,
P_Z(d(x,b,y)).
\]

For a deterministic input sample
\[
\mathbf z=((x_i,b_i,y_i))_{i=1}^n
\in\bigl(\{1,\ldots,d\}\times\{0,1\}^2\bigr)^n
\]
and \(\mathbf u=(u_i)_{i=1}^n\in\{0,1\}^n\), define
\[
\mathsf o_n(\mathbf z,\mathbf u)
:=
\bigl(o(x_i,b_i,y_i;u_i)\bigr)_{i=1}^n,
\qquad
\kappa_{\mathrm{syn},n}(\mathbf z,\cdot)
:=
\bigotimes_{i=1}^n
\kappa_{\mathrm{syn}}\bigl((x_i,b_i,y_i),\cdot\bigr).
\]
Independent fair draws give
\[
\kappa_{\mathrm{syn},n}(\mathbf z,\cdot)
=
\bigl(\mathsf o_n(\mathbf z,\cdot)\bigr)_{\#}
\nu_{\mathrm{coin}}^{\otimes n},
\]
where the subscript \(\#\) denotes pushforward. Integrating over the input sample, regrouping the product coordinates into record--coin pairs, and then applying the record map coordinatewise gives
\[
\begin{aligned}
\int\kappa_{\mathrm{syn},n}(\mathbf z,\cdot)\,
P_Z^{\otimes n}(d\mathbf z)
&=
(\mathsf o_n)_{\#}
\bigl(P_Z^{\otimes n}\otimes\nu_{\mathrm{coin}}^{\otimes n}\bigr)\\
&=
\bigotimes_{i=1}^n
\left[
\int\kappa_{\mathrm{syn}}\bigl((x,b,y),\cdot\bigr)\,
P_Z(d(x,b,y))
\right]\\
&=
\nu_{\mathrm{obs}}^{\otimes n}.
\end{aligned}
\]
The regrouping is justified by the product identity
\[
\left(\prod_{i=1}^nP_Z(dz_i)\right)
\left(\prod_{i=1}^n\nu_{\mathrm{coin}}(du_i)\right)
=
\prod_{i=1}^n
\left(P_Z(dz_i)\nu_{\mathrm{coin}}(du_i)\right),
\]
and coordinatewise pushforward preserves this product form. Thus the kernel produces exactly the observed product law identified in the preceding step, and its output satisfies \cref{obj:ass:iid-sampling}.
% lean: CausalSmith.Stat.MarRareqLogfrontier.SyntheticCompletion.selected_eq
% lean: CausalSmith.Stat.MarRareqLogfrontier.SyntheticCompletion.observedLaw_eq_bind
% lean: CausalSmith.Stat.MarRareqLogfrontier.syntheticProductLaw_eq_coin_map
% lean: CausalSmith.Stat.MarRareqLogfrontier.syntheticSample_bind_eq_pi

\item Compute the target through cell contributions. Since \(P_Z^*\) belongs to the arrival-floor model, \cref{obj:lem:unrestricted-cell-identification} supplies
\[
\tau(P_Z^*)=\Phi(P_Z^*)
=
2\sum_{(a,x,s)\in\mathcal J_d}g(a)c_{axs}(P_Z^*).
\]
Here the contributions and their zero convention are given by \cref{obj:def:cell-functional}, and \cref{obj:synth_19} gives \(g(0)=-1\), \(g(1)=1\).

For \(p_x>0\), the arrival mass in either arm is strictly positive:
\[
P_Z^*(A=a,X=x-1,S=0,R=1)
=
\frac{p_x}{2}\beta^{\mathrm{mark}}_{ax}
\ge\frac{qp_x}{2}>0.
\]
The joint-mass identity in \cref{eq:synthetic-randomization-joint-mass} therefore yields
\[
\mu_{a,x-1,0}(P_Z^*)
=
\frac{\frac{p_x}{2}\beta^{\mathrm{mark}}_{ax}\mu_{ax}}
     {\frac{p_x}{2}\beta^{\mathrm{mark}}_{ax}}
=\mu_{ax},
\qquad
c_{a,x-1,0}(P_Z^*)=\frac{p_x}{2}\mu_{ax}.
\]
When \(p_x=0\), the arrived-cell probability is zero, so the defining convention gives the same contribution identity. Every cell with \(S=1\) also has zero arrived-cell probability. Hence, for every category and arm,
\[
c_{a,x-1,0}(P_Z^*)=\frac{p_x}{2}\mu_{ax},
\qquad
c_{a,x-1,1}(P_Z^*)=0.
\]
Substituting these identities into the signed cell sum gives
\[
\begin{aligned}
\tau(P_Z^*)
&=2\sum_{x=1}^d
\left(\frac{p_x}{2}\mu_{1x}
-\frac{p_x}{2}\mu_{0x}\right)\\
&=\sum_{x:p_x>0}p_x(\mu_{1x}-\mu_{0x})
=\theta_Z(P_Z),
\end{aligned}
\]
where the last equality uses \cref{obj:def:zeng-ate-functional}; null categories contribute zero.
% lean: CausalSmith.Stat.MarRareqLogfrontier.SyntheticCompletion.full_cellContribution_false
% lean: CausalSmith.Stat.MarRareqLogfrontier.SyntheticCompletion.full_cellContribution_true
% lean: CausalSmith.Stat.MarRareqLogfrontier.SyntheticCompletion.fullLaw_ate
\end{enumerate}
\end{proof}
